% Derivation of the exact solution starting in the frequency domain.
% Draft of 2026-11-21. Companion to derivation-exact-solution.tex (time domain).
% Equation numbers of the theory note refer to its version of 2026-11-21.
\documentclass[aps,pra,onecolumn,notitlepage,11pt,amsmath,amssymb,nofootinbib]{revtex4-2}

\usepackage{bm}
\usepackage{array}
\usepackage[colorlinks=true,allcolors=blue]{hyperref}
\newcolumntype{R}[1]{>{\raggedright\arraybackslash}p{#1}}

\newcommand{\Nop}{\hat N}
\newcommand{\Iop}{\hat I}
\newcommand{\dd}{\mathrm{d}}
\newcommand{\ee}{\mathrm{e}}
\newcommand{\ii}{\mathrm{i}}
\newcommand{\sinc}{\operatorname{sinc}}
\newcommand{\PV}{\mathcal P}
\newcommand{\TN}[1]{\text{TN\,#1}}

\begin{document}

\title{Optical rectification with a quantized pump:\\
derivation starting in the frequency domain}
\author{Working notes, 21 November 2026}
\noaffiliation

\begin{abstract}
The exact solution of the theory note is derived here with the mode
operators $\hat a(\omega)$ and $\hat b(\Omega)$ from the start. In this form
the Hamiltonian is the usual three-wave-mixing Hamiltonian with a coupling
function, the condition for solvability is a one-line statement about that
function, the equation of motion of each THz mode is that of a driven
oscillator, and the solution is the coupling on the energy shell. The
evolution operator separates into an on-shell part, which is the
displacement of the THz field, and an off-shell part, which is a phase given
by a principal-value integral. The one step that is easier in the time
domain is the transformation of the pump, where a convolution in frequency
becomes a multiplication by a phase. Tags (TN\,8)--(TN\,19) are equation
numbers of the theory note of 21 November 2026, whose Sec.~II follows the
order of these notes; Eqs.~\eqref{eq:Hgen}, \eqref{eq:Gamma}, \eqref{eq:g},
\eqref{eq:N}, and \eqref{eq:Hint} below are its Eqs.~(1), (2), (4), (5), and
(6). The time-domain derivation is in
\texttt{derivation-exact-solution.pdf}.
\end{abstract}

\maketitle

\section{Summary}

\begin{center}
\begin{tabular}{R{0.30\textwidth}R{0.31\textwidth}R{0.31\textwidth}}
\hline\hline
Step & Frequency domain & Time domain\\
\hline
Hamiltonian & three-wave mixing with a coupling function $g(\Omega,\Omega')$, Eq.~\eqref{eq:Hint} & flux at a retarded time times the THz field along the crystal\\[4pt]
Condition for solvability & $g$ independent of the optical frequency; $[\Nop,\Nop^\dagger]=0$ by a shift of the integration variable, Eq.~\eqref{eq:F2} & $[\Iop(\tau),\Iop(\tau')]=0$\\[4pt]
THz field & driven oscillator; $\zeta=2\pi g$ on the energy shell, Eq.~\eqref{eq:onshell} & integral of a source along $z$ and $t$\\[4pt]
Pump & ladder of coupled sidebands, a convolution, Eq.~\eqref{eq:ladder} & multiplication by a phase: simpler\\[4pt]
Evolution operator & on-shell displacement and off-shell phase, Eqs.~\eqref{eq:D}, \eqref{eq:PhiPV} & kernel $W(\tau-\tau')$ with a sign function\\[4pt]
State of the THz field & not transparent & mixture over photon arrival times: simpler\\
\hline\hline
\end{tabular}
\end{center}

\section{Hamiltonian}
\label{sec:H}

\subsection{Operators}

Schr\"odinger picture. The pump has mode operators $\hat a(\omega)$ and the
THz field $\hat b(\Omega)$, $\Omega>0$,
\begin{equation}
[\hat a(\omega),\hat a^\dagger(\omega')]=\delta(\omega-\omega'),\qquad
[\hat b(\Omega),\hat b^\dagger(\Omega')]=\delta(\Omega-\Omega'),
\end{equation}
and pump operators commute with THz operators. The fields in the crystal are
\begin{equation}
\hat E_o^{(+)}(z)=\ii\sqrt{\frac{\hbar\omega_0}{4\pi\varepsilon_0cn_oA}}
\int\!\dd\omega\;\hat a(\omega)\,\ee^{\ii k(\omega)z},\qquad
\hat E_T(z)=-\int_0^\infty\!\!\dd\Omega\;s(\Omega)\Big[\hat b(\Omega)\,\ee^{\ii k_T(\Omega)z}+\mathrm{H.c.}\Big],
\label{eq:fields}
\end{equation}
with $s(\Omega)=[\hbar\Omega/(4\pi\varepsilon_0cn_TA)]^{1/2}$, and the free
Hamiltonian is
\begin{equation}
\hat H_0=\hbar\int_0^\infty\!\!\dd\Omega\;\Omega\,\hat b^\dagger(\Omega)\hat b(\Omega)
+\hbar\int\!\dd\omega\;\omega\,\hat a^\dagger(\omega)\hat a(\omega).
\label{eq:H0}
\end{equation}

\subsection{Interaction}

The rectified polarization
$\hat P_{\rm OR}=2\varepsilon_0\chi^{(2)}\hat E_o^{(-)}\hat E_o^{(+)}$ contains
all pairs of pump frequencies, and the interaction energy
$-A\int_0^L\dd z\,\hat P_{\rm OR}\hat E_T$ is, with Eq.~\eqref{eq:fields},
\begin{equation}
\hat H_{\rm int}=\hbar\int_0^\infty\!\!\dd\Omega\iint\!\dd\omega\,\dd\omega'\;
\Big[\Gamma(\Omega;\omega,\omega')\,\hat b^\dagger(\Omega)\,\hat a^\dagger(\omega)\,\hat a(\omega')
+\mathrm{H.c.}\Big],
\label{eq:Hgen}
\end{equation}
\begin{equation}
\Gamma(\Omega;\omega,\omega')=\frac{\kappa(\Omega)}{2\pi}\int_0^L\!\!\dd z\,\ee^{-\ii\Delta z}
=\frac{\kappa(\Omega)L}{2\pi}\,\ee^{-\ii\Delta L/2}\,\sinc\!\Big(\frac{\Delta L}{2}\Big),
\qquad \Delta=k_T(\Omega)+k(\omega)-k(\omega'),
\label{eq:Gamma}
\end{equation}
with $\kappa(\Omega)=(\chi^{(2)}\omega_0/n_oc)\,s(\Omega)$ as in the theory note.
This is the three-wave-mixing Hamiltonian in its usual form. Each term
annihilates a pump photon at $\omega'$ and creates a pump photon at $\omega$
and a THz photon at $\Omega$; $\Delta$ is the wave-number mismatch of that
process, and the crystal of length $L$ allows $|\Delta|\lesssim2\pi/L$. The
Hamiltonian does not impose $\omega'=\omega+\Omega$. Energy conservation is
produced by the time evolution, Sec.~\ref{sec:thz}.

\subsection{Flat coupling}

Without group-velocity dispersion of the pump,
$k(\omega')-k(\omega)=(\omega'-\omega)/v_g$, and $\chi^{(2)}$ independent of the
optical frequency, the coupling function depends on the two optical
frequencies only through their difference $\Omega'=\omega'-\omega$:
\begin{equation}
\Gamma(\Omega;\omega,\omega+\Omega')=g(\Omega,\Omega'),\qquad
g(\Omega,\Omega')=\frac{\kappa(\Omega)L}{2\pi}\,\ee^{-\ii\delta L/2}\,\sinc\!\Big(\frac{\delta L}{2}\Big),
\qquad\delta(\Omega,\Omega')=k_T(\Omega)-\frac{\Omega'}{v_g}.
\label{eq:g}
\end{equation}
\emph{This is the assumption of the theory, stated in one line.} The double
integral over the optical frequencies then contains the pump only through
\begin{equation}
\Nop(\Omega')=\int\!\dd\omega\;\hat a^\dagger(\omega)\,\hat a(\omega+\Omega'),\qquad
\Nop^\dagger(\Omega')=\Nop(-\Omega'),
\label{eq:N}
\end{equation}
and
\begin{equation}
\boxed{\;\hat H_{\rm int}=\hbar\int_0^\infty\!\!\dd\Omega\int_{-\infty}^{\infty}\!\!\dd\Omega'\;
\Big[g(\Omega,\Omega')\,\hat b^\dagger(\Omega)\,\Nop(\Omega')+g^*(\Omega,\Omega')\,\hat b(\Omega)\,\Nop^\dagger(\Omega')\Big]\;}
\label{eq:Hint}
\end{equation}
$\Nop(\Omega')$ is the autocorrelation of the pump spectrum at the lag
$\Omega'$. In classical nonlinear optics the source of rectification at
$\Omega$ is $\int\dd\omega\,E^*(\omega)E(\omega+\Omega)$; Eq.~\eqref{eq:N} is the
same expression with operators, normally ordered. It is the Fourier transform
of the photon flux, $\Nop(\Omega)=\int\dd\tau\,\ee^{\ii\Omega\tau}\Iop(\tau)$,
and $\Nop(0)$ is the photon number.

\section{Algebra of the source operator}
\label{sec:alg}

\emph{(F1) $\Nop$ shifts the pump frequency.} From the commutator of the mode
operators,
\begin{equation}
[\Nop(\Omega),\hat a(\omega)]=-\hat a(\omega+\Omega),\qquad
[\Nop(\Omega),\hat a^\dagger(\omega)]=\hat a^\dagger(\omega-\Omega).
\label{eq:F1}
\end{equation}

\emph{(F2) The source operators commute.} With
$[\hat a^\dagger(\omega)\hat a(\omega+\Omega),\hat a^\dagger(\omega_1)\hat a(\omega_1+\Omega_1)]
=\delta(\omega+\Omega-\omega_1)\,\hat a^\dagger(\omega)\hat a(\omega_1+\Omega_1)
-\delta(\omega_1+\Omega_1-\omega)\,\hat a^\dagger(\omega_1)\hat a(\omega+\Omega)$,
\begin{equation}
[\Nop(\Omega),\Nop(\Omega_1)]
=\int\!\dd\omega\;\hat a^\dagger(\omega)\,\hat a(\omega+\Omega+\Omega_1)
-\int\!\dd\omega_1\;\hat a^\dagger(\omega_1)\,\hat a(\omega_1+\Omega_1+\Omega)=0 ,
\label{eq:F2}
\end{equation}
for all real $\Omega$, $\Omega_1$, hence also
$[\Nop(\Omega),\Nop^\dagger(\Omega_1)]=0$. The two integrals are identical.
Two things were needed. First, the weight of every pair
$(\omega,\omega+\Omega)$ in $\Nop$ is the same, namely one; this is the flat
coupling of Eq.~\eqref{eq:g}. With a weight that depends on $\omega$ the two
integrals would differ. Second, the integration range has no edge; with a
band of finite width the integrals differ by the population at the band
edges. Both statements are immediate in the frequency domain, and they are
the complete condition for what follows.

\emph{(F3) Free evolution.} From Eq.~\eqref{eq:H0},
$[\hat H_0,\hat a^\dagger(\omega)\hat a(\omega+\Omega)]
=\hbar[\omega-(\omega+\Omega)]\,\hat a^\dagger(\omega)\hat a(\omega+\Omega)$, so
\begin{equation}
\frac{\ii}{\hbar}[\hat H_0,\Nop(\Omega)]=-\ii\Omega\,\Nop(\Omega),\qquad
\frac{\ii}{\hbar}[\hat H_0,\hat b(\Omega)]=-\ii\Omega\,\hat b(\Omega),\qquad
\frac{\ii}{\hbar}[\hat H_0,\hat a(\omega)]=-\ii\omega\,\hat a(\omega).
\label{eq:F3}
\end{equation}
The beat note of the pump at $\Omega$ oscillates at $\Omega$, like a THz mode.

\section{Heisenberg equations and the THz field: Eqs.~(8)--(11)}
\label{sec:thz}

Heisenberg operators evolve from the remote past, when the pulse has not yet
reached the crystal, with
$\partial_t\hat O_H=(\ii/\hbar)[\hat H_0+\hat H_{\rm int},\hat O]_H$.

\emph{Source.} By (F2), $[\hat H_{\rm int},\Nop(\Omega')]=0$, and with (F3)
\begin{equation}
\frac{\partial\Nop_H(\Omega',t)}{\partial t}=-\ii\Omega'\,\Nop_H(\Omega',t)
\quad\Longrightarrow\quad
\Nop_H(\Omega',t)=\Nop(\Omega')\,\ee^{-\ii\Omega't}.
\label{eq:Nsol}
\end{equation}
Every Fourier component of the pump intensity keeps its input amplitude and
oscillates freely. This is the conservation of the photon flux,
\begin{equation}
\boxed{\;\Nop_{\rm out}(\Omega)=\Nop_{\rm in}(\Omega),\qquad\text{equivalently}\qquad
\Iop_{\rm out}(\tau)=\Iop_{\rm in}(\tau)\;}
\tag{\TN{8}}
\end{equation}

\emph{THz field.} From Eq.~\eqref{eq:Hint},
$(\ii/\hbar)[\hat H_{\rm int},\hat b(\Omega)]
=-\ii\int\dd\Omega'\,g(\Omega,\Omega')\Nop(\Omega')$, so with
Eq.~\eqref{eq:Nsol}
\begin{equation}
\frac{\partial\hat b_H(\Omega,t)}{\partial t}=-\ii\Omega\,\hat b_H(\Omega,t)
-\ii\int_{-\infty}^{\infty}\!\!\dd\Omega'\;g(\Omega,\Omega')\,\Nop(\Omega')\,\ee^{-\ii\Omega't}.
\label{eq:bH}
\end{equation}
Each THz mode is a harmonic oscillator of frequency $\Omega$ driven by the
beat notes of the pump, with the coupling function as the weight. The source
is the input operator. Removing the free oscillation,
$\hat b(\Omega,t)=\ee^{\ii\Omega t}\hat b_H(\Omega,t)$,
\begin{equation}
\boxed{\;\frac{\partial\hat b(\Omega,t)}{\partial t}
=-\ii\int_{-\infty}^{\infty}\!\!\dd\Omega'\;g(\Omega,\Omega')\,\Nop(\Omega')\,\ee^{\ii(\Omega-\Omega')t}\;}
\tag{\TN{9}}
\end{equation}
Integrating over all times, the oscillating factor gives
$\int\dd t\,\ee^{\ii(\Omega-\Omega')t}=2\pi\delta(\Omega-\Omega')$: only the
beat note at the frequency of the oscillator drives it. Hence
\begin{equation}
\boxed{\;\hat b_{\rm out}(\Omega)=\hat b_{\rm in}(\Omega)-\ii\,\zeta(\Omega)\,\Nop(\Omega)\;}
\tag{\TN{10}}
\end{equation}
\begin{equation}
\boxed{\;\zeta(\Omega)=2\pi\,g(\Omega,\Omega)
=\kappa(\Omega)L\,\ee^{-\ii\Delta kL/2}\,\sinc\!\Big(\frac{\Delta kL}{2}\Big),
\qquad\Delta k(\Omega)=\delta(\Omega,\Omega)=k_T(\Omega)-\frac{\Omega}{v_g}\;}
\tag{\TN{11}}
\label{eq:onshell}
\end{equation}
The amplitude $\zeta$ is $2\pi$ times the coupling function on the energy
shell. The derivation has three lines, and no step is an approximation.

\emph{Comparison with a general coupling.} For the Hamiltonian
\eqref{eq:Hgen} with a coupling that depends on the optical frequency, the
same steps to first order in $\Gamma$ give
$\hat b_{\rm out}=\hat b_{\rm in}
-2\pi\ii\int\dd\omega\,\Gamma(\Omega;\omega,\omega+\Omega)\,
\hat a^\dagger(\omega)\hat a(\omega+\Omega)$. That is the result of perturbation
theory, and it is not exact, because its source does not commute with the
Hamiltonian. With flat coupling the source is $\Nop(\Omega)$, a constant of
the motion, and first order is the whole answer.

\section{Pump: Eqs.~(14)--(17)}
\label{sec:pump}

With (F1), the Heisenberg equation of a pump mode is
\begin{align}
\frac{\partial\hat a_H(\omega,t)}{\partial t}=-\ii\omega\,\hat a_H(\omega,t)
-\ii\int_0^\infty\!\!\dd\Omega\int_{-\infty}^{\infty}\!\!\dd\Omega'\,
\Big[&g(\Omega,\Omega')\,\hat b_H^\dagger(\Omega,t)\,\hat a_H(\omega+\Omega',t)\nonumber\\
&+g^*(\Omega,\Omega')\,\hat b_H(\Omega,t)\,\hat a_H(\omega-\Omega',t)\Big].
\label{eq:ladder}
\end{align}
The amplitude at $\omega$ is fed from $\omega+\Omega'$ with the emission of a
THz photon and from $\omega-\Omega'$ with the absorption of one. The modes
form a ladder of coupled sidebands; this is cascading. In the frequency
domain the equation is a convolution and is not solved by inspection.

A convolution in frequency is a product in time. With the envelope operator
\begin{equation}
\hat a(\tau)=\frac{1}{\sqrt{2\pi}}\int\!\dd\omega\;\hat a(\omega)\,\ee^{-\ii(\omega-\omega_0)\tau},
\qquad
\frac{1}{\sqrt{2\pi}}\int\!\dd\omega\;\hat a(\omega\pm\Omega')\,\ee^{-\ii(\omega-\omega_0)\tau}
=\ee^{\pm\ii\Omega'\tau}\,\hat a(\tau),
\end{equation}
and
$\int\dd\Omega'\,g(\Omega,\Omega')\,\ee^{\ii\Omega'\tau}
=\kappa(\Omega)\int_0^L\dd z\,\ee^{-\ii k_Tz}\,\delta(\tau+z/v_g)$, which is
non-zero only for a slice inside the crystal, Eq.~\eqref{eq:ladder} becomes
\begin{equation}
\frac{\partial\hat a_H(\tau,t)}{\partial t}
=\Big(\frac{\partial}{\partial\tau}-\ii\omega_0\Big)\hat a_H(\tau,t)
-\ii v_g\,\hat\Theta_H(-v_g\tau,t)\,\hat a_H(\tau,t),
\end{equation}
for $0<-v_g\tau<L$, with
$\hat\Theta_H(z,t)=\int\dd\Omega\,\kappa\,[\hat b_H(\Omega,t)\,\ee^{\ii k_Tz}+\mathrm{H.c.}]$.
The envelope of each slice is multiplied by a phase
factor at a rate set by the THz field at its position. This is (TN\,14), and
from here the time-domain derivation applies without change
(\texttt{derivation-exact-solution.pdf}, Sec.~VIII):
\begin{equation}
\hat a_{\rm out}(\tau)=\ee^{-\ii\hat\vartheta(\tau)}\,\hat a_{\rm in}(\tau),\qquad
\hat\vartheta(\tau)=\tfrac12\big[\hat F_{\rm in}(\tau)+\hat F_{\rm out}(\tau)\big]+\hat\theta(\tau),\qquad
\hat F(\tau)=\int\!\dd\Omega\,\big[\zeta\,\hat b^\dagger\ee^{\ii\Omega\tau}+\mathrm{H.c.}\big].
\tag{\TN{15}--\TN{17}}
\end{equation}
Back in the frequency domain the result is a convolution with an
operator-valued kernel,
\begin{equation}
\hat a_{\rm out}(\omega)=\int\!\dd\omega'\;\hat{\mathcal M}(\omega-\omega')\,\hat a_{\rm in}(\omega'),
\qquad
\hat{\mathcal M}(\nu)=\frac{1}{2\pi}\int\!\dd\tau\;\ee^{\ii\nu\tau}\,\ee^{-\ii\hat\vartheta(\tau)} .
\end{equation}
For a classical monochromatic THz field, $\vartheta(\tau)=m\cos(\Omega\tau+\phi)$,
the kernel is the series of sidebands at $\nu=n\Omega$ with Bessel-function
amplitudes $J_n(m)$, the familiar spectrum of a phase-modulated wave. The
ladder equation \eqref{eq:ladder} generates this series order by order; the
time-domain form gives it at once.

\section{Evolution operator: Eqs.~(18) and (19)}
\label{sec:U}

In the interaction picture with respect to $\hat H_0$, by (F3),
\begin{equation}
\hat H_{\rm int}(t)=\hbar\int_0^\infty\!\!\dd\Omega\int_{-\infty}^{\infty}\!\!\dd\Omega'\,
\Big[g(\Omega,\Omega')\,\hat b^\dagger(\Omega)\Nop(\Omega')\,\ee^{\ii(\Omega-\Omega')t}+\mathrm{H.c.}\Big],
\label{eq:HI}
\end{equation}
and $\hat U=\mathcal T\exp[-(\ii/\hbar)\int\dd t\,\hat H_{\rm int}(t)]$. The
commutator of $\hat H_{\rm int}$ at two times contains $\Nop$ and
$\Nop^\dagger$ only, by $[\hat b^\dagger(\Omega)\Nop(\Omega'),\hat b(\Omega_1)\Nop^\dagger(\Omega'')]
=-\delta(\Omega-\Omega_1)\Nop(\Omega')\Nop^\dagger(\Omega'')$ and (F2), and it
commutes with $\hat H_{\rm int}$. The Magnus expansion therefore ends at its
second term,
$\hat U=\ee^{\hat\Omega_2}\ee^{\hat\Omega_1}$ (verified by differentiation in
\texttt{derivation-exact-solution.pdf}, Sec.~IX\,B).

\emph{First term: on the energy shell.} The time integral of
Eq.~\eqref{eq:HI} gives $2\pi\delta(\Omega-\Omega')$:
\begin{equation}
\hat\Omega_1=-\frac{\ii}{\hbar}\int\!\dd t\,\hat H_{\rm int}(t)
=-\ii\int_0^\infty\!\!\dd\Omega\,\Big[2\pi g(\Omega,\Omega)\,\hat b^\dagger(\Omega)\Nop(\Omega)+\mathrm{H.c.}\Big],
\end{equation}
\begin{equation}
\boxed{\;\hat D=\ee^{\hat\Omega_1}=\exp\!\Big\{-\ii\int_0^\infty\!\!\dd\Omega\,
\big[\zeta(\Omega)\,\hat b^\dagger(\Omega)\Nop(\Omega)+\mathrm{H.c.}\big]\Big\}\;}
\tag{\TN{19}}
\label{eq:D}
\end{equation}

\emph{Second term: off the energy shell.} With
$\int\dd t\int^t\dd t'\,\ee^{\ii\alpha t-\ii\beta t'}
=2\pi\delta(\alpha-\beta)\,[\pi\delta(\beta)+\ii\PV(1/\beta)]$ for the two
orderings of the commutator, the delta-function parts cancel and the
principal values add:
\begin{equation}
\hat\Omega_2=-\frac{1}{2\hbar^2}\int\!\dd t\!\int^t\!\!\dd t'\,
[\hat H_{\rm int}(t),\hat H_{\rm int}(t')]=\ii\hat\Phi,
\end{equation}
\begin{equation}
\boxed{\;\hat\Phi=\int_{-\infty}^{\infty}\!\!\dd\Omega'\;\Nop^\dagger(\Omega')\Nop(\Omega')\;
2\pi\,\PV\!\!\int_0^\infty\!\!\dd\Omega\,\frac{|g(\Omega,\Omega')|^2}{\Omega-\Omega'}\;}
\label{eq:PhiPV}
\end{equation}
and
\begin{equation}
\boxed{\;\hat U=\ee^{\ii\hat\Phi}\,\hat D\;}
\tag{\TN{18}}
\end{equation}
The structure is that of scattering theory. The coupling on the energy shell
produces real emission: the displacement $\hat D$. The coupling off the shell
produces no emission but a phase that depends on the flux, through the
virtual emission and reabsorption of THz photons by the pump: a principal
value integral over the squared coupling, as for a level shift. In optics
this phase is the cascaded Kerr effect of rectification and the
electro-optic effect.

\emph{Evaluation.} The coupling function $|g(\Omega,\Omega')|^2$ is a
$\sinc^2$ of width $2\pi/L$ in the wave number $\delta$, centred where
$k_T(\Omega)=\Omega'/v_g$. If $\kappa^2$ is taken at the shell, constant over
that width, and $k_T$ is linear over it, the integral is elementary, with
the triangle function as the Fourier transform of $\sinc^2$:
\begin{equation}
\PV\!\!\int_{-\infty}^{\infty}\!\!\dd u\;\frac{1}{u}\,\sinc^2\!\Big(\frac{(u+\Delta k)L}{2}\Big)
=-\frac{2\pi}{L^2}\,\frac{\Delta kL-\sin\Delta kL}{\Delta k^2},
\label{eq:pvint}
\end{equation}
\begin{equation}
\hat\Phi\simeq-\int_0^\infty\!\!\dd\Omega\;\varphi(\Omega)\,\Nop^\dagger(\Omega)\Nop(\Omega),
\qquad
\varphi(\Omega)=\kappa^2\,\frac{\Delta kL-\sin\Delta kL}{\Delta k^2}.
\label{eq:phi}
\end{equation}
This is Eq.~(A7) of the theory note. It is also what ordering the propagation
in $z$ gives, with the generator of Eq.~(A8) there, which contains the
coupling on the shell only. \emph{The propagation picture is therefore the
approximation in which the coupling strength is evaluated on the energy
shell.} The variation of $\kappa^2\propto\Omega$ across the width of the
coupling function gives the correction: writing
$\Omega=\Omega'+(\Omega-\Omega')$ in the numerator of Eq.~\eqref{eq:PhiPV},
the second part has no pole and integrates to a constant,
\begin{equation}
\hat\Phi=-\int_0^\infty\!\!\dd\Omega\,\Big[\varphi(\Omega)-c_\kappa\frac{cL}{n_T}\Big]\Nop^\dagger\Nop
\qquad(\kappa^2=c_\kappa\Omega,\ \text{dispersionless THz index}),
\label{eq:Kerr}
\end{equation}
and $\int_0^\infty\dd\Omega\,\Nop^\dagger\Nop=\pi\int\dd\tau\,\Iop(\tau)^2$: an
instantaneous Kerr-type phase, linear in $L$. It is the term by which the
time-ordered kernel $W$ differs from the kernel $V$ in the time-domain
notes. The terms of Eq.~\eqref{eq:Hint} with $\Omega'<0$, which do not
conserve energy, contribute to $\hat\Phi$ an amount that does not grow with
$L$.

\emph{Backward waves.} The model contains forward THz waves only. For
backward waves the coupling function has $\delta_b=-k_T(\Omega)-\Omega'/v_g$,
which vanishes only for $\Omega'<0$: the wave numbers match for a process
that does not conserve energy. There is then no pole, and the integral in
Eq.~\eqref{eq:PhiPV} is the area of the peak of $|g_b|^2$ divided by its
distance $\Omega+|\Omega'|$ from the shell. For a dispersionless THz index and
$\kappa^2=c_\kappa\Omega$ this gives the coefficient
$c_\kappa(cL/n_T)\,n_g/(n_g+n_T)$ in front of $\Nop^\dagger\Nop$, with
$n_g=c/v_g$: a term of the same form as the constant in Eq.~\eqref{eq:Kerr}
and about half its size. This is an estimate and has not been checked
numerically. It shows that the instantaneous Kerr-type phase is not
determined by a model with forward waves only.

\section{Relation to the coupled-wave equations of nonlinear optics}

If the evolution is ordered in $z$ and the coupling is taken on the shell,
the equations of motion are
\begin{equation}
\frac{\partial\hat b(\Omega,z)}{\partial z}=-\ii\kappa(\Omega)\,\ee^{-\ii\Delta kz}\,\Nop(\Omega),\qquad
\frac{\partial\Nop(\Omega)}{\partial z}=0,
\end{equation}
\begin{equation}
\frac{\partial\hat a(\omega,z)}{\partial z}=-\ii\int_0^\infty\!\!\dd\Omega\;\kappa(\Omega)
\Big[\ee^{-\ii\Delta kz}\,\hat b^\dagger(\Omega,z)\,\hat a(\omega+\Omega,z)
+\ee^{\ii\Delta kz}\,\hat b(\Omega,z)\,\hat a(\omega-\Omega,z)\Big].
\end{equation}
The first is the textbook equation for the growth of the THz amplitude, with
operators:
\begin{equation}
\frac{\partial E_T(\Omega,z)}{\partial z}\;\propto\;-\ii\,\ee^{-\ii\Delta kz}
\int\!\dd\omega\;E^*(\omega)\,E(\omega+\Omega)\qquad\text{(classical)}.
\end{equation} The classical theory solves it in the undepleted-pump
approximation by holding the source fixed. The quantum statement is that the
source is fixed exactly, although the pump spectrum $\hat a(\omega,z)$ itself
changes along the crystal by the second equation: the spectrum is reshaped,
its autocorrelation is not.

\section{Checks}

\begin{center}
\begin{tabular}{R{0.38\textwidth}R{0.55\textwidth}}
\hline\hline
Statement & Check\\
\hline
Principal-value integral, Eq.~\eqref{eq:pvint} & numerical integration, four parameter sets: agreement to six digits; two of them are in the test suite\\[4pt]
Eq.~\eqref{eq:Kerr}: full coefficient of Eq.~\eqref{eq:PhiPV} for $\kappa^2=c_\kappa\Omega$, dispersionless index & numerical integration, five parameter sets ($L\Omega n_T/c$ from 20 to 200): $-\varphi+c_\kappa cL/n_T$ reproduced to better than $10^{-4}$ relative, after adding the terms with $\Omega'<0$\\[4pt]
Terms with $\Omega'<0$ & same integration: of order $0.16\,c_\kappa$ at $\Omega'=1$ in these units, independent of $L$\\[4pt]
(TN\,8), (TN\,10), (TN\,11), (TN\,18), (TN\,19) & the tests in \texttt{code/tests} for optical rectification; their model, a single generator $g(\hat b^\dagger\Nop+\hat b\Nop^\dagger)$ with $\Nop=\sum_j\hat a_j^\dagger\hat a_{j+1}$, is Eq.~\eqref{eq:Hint} with one THz mode on a discrete frequency ladder\\
\hline\hline
\end{tabular}
\end{center}

One case of each numerical integration is in the test suite (the test on the
off-shell phase in the same file). Not checked: Eq.~\eqref{eq:PhiPV} with THz
dispersion or absorption.

\section{Which route for which purpose}

For the main line of the theory note, Hamiltonian, solvability, and the THz
output, the frequency domain is shorter and each step has a counterpart in
the classical theory of rectification and in the theory of parametric
down-conversion. For the pump and for the state of the THz field the time
domain is simpler, because the conserved quantity is local in time: the flux
$\Iop(\tau)$ is diagonal in the basis of photon arrival times, and the action
of the THz field on the pump is a multiplication. A presentation that starts
in the frequency domain and changes to the time domain after (TN\,11) uses
each where it is shortest. Section~II of the theory note of 21 November 2026
is written in this order.

\end{document}
